\DocumentMetadata{
  pdfversion=2.0,
  pdfstandard=ua-2,
  lang=en-US,
  tagging=on,
  tagging-setup={math/setup=mathml-SE,tabsorder=structure}
}
\documentclass[11pt,letterpaper]{article}
\usepackage[margin=0.68in,includefoot]{geometry}
\usepackage{newcomputermodern}
\usepackage{amsmath,amscd,mathtools}
\usepackage{tikz,adjustbox}
\providecommand{\square}{\mdlgwhtsquare}
\usepackage[hidelinks]{hyperref}
\hypersetup{
  pdftitle={Numerical Analysis PhD Exam, May 2010},
  pdfauthor={University of Florida Department of Mathematics},
  pdfsubject={Graduate mathematics examination},
  pdfdisplaydoctitle=true,
  bookmarksopen=true
}
\setcounter{secnumdepth}{0}
\setlength{\parindent}{0pt}
\setlength{\parskip}{0.28em}
\setlength{\emergencystretch}{3em}
\setlength{\leftmargini}{2.15em}
\setlength{\leftmarginii}{2.65em}
\setlength{\leftmarginiii}{3em}
\renewcommand{\labelenumi}{\textbf{\theenumi.}}
\pagestyle{empty}
\raggedbottom
\allowdisplaybreaks
\newenvironment{examproblems}[1]{%
  \begin{enumerate}%
  \setcounter{enumi}{#1}%
  \setlength{\topsep}{0.35em}%
  \setlength{\partopsep}{0pt}%
  \setlength{\itemsep}{0.48em}%
  \setlength{\parsep}{0.18em}%
}{\end{enumerate}}
\newenvironment{examsubparts}{%
  \begin{enumerate}%
  \setlength{\topsep}{0.18em}%
  \setlength{\partopsep}{0pt}%
  \setlength{\itemsep}{0.16em}%
  \setlength{\parsep}{0.1em}%
}{\end{enumerate}}
\newenvironment{examinstructions}{%
  \par\begingroup\itshape
}{%
  \par\endgroup\vspace{0.18em}
}
\makeatletter
\renewcommand\section{\@startsection{section}{1}{0pt}%
  {-1.25ex plus -.25ex minus -.1ex}%
  {0.55ex plus .1ex}%
  {\normalfont\large\bfseries}}
\renewcommand{\@maketitle}{%
  \newpage\null\vskip -1em%
  {\footnotesize\bfseries
    \href{https://www.ufl.edu/}{University of Florida}, \href{https://math.ufl.edu/}{Department of Mathematics}\par}%
  \vskip -0.5em%
  \tagstructbegin{tag=Title}%
    {\LARGE\bfseries\href{https://gma.math.ufl.edu/exams/numerical-analysis-phd/}{Numerical Analysis PhD Exam}, May 2010\par}%
  \tagstructend%
  \vskip 0.25em%
  \tagmcbegin{artifact}%
    \hrule height 1pt%
  \tagmcend%
  \vskip 1em%
}
\makeatother
\title{Numerical Analysis PhD Exam, May 2010}
\author{}
\date{}
\begin{document}
\maketitle
\thispagestyle{empty}
\begin{examinstructions}
Answer any 8 questions.
\end{examinstructions}
\begin{examproblems}{0}
\item
If~\(P\) is a projector that is neither~\(0\) nor the identity, prove that \(\lVert P\rVert=\lVert I-P\rVert\) in any norm induced by a vector norm generated by an inner product.
\item
Let \(M=\begin{bmatrix}A&B^t\\B&C\end{bmatrix}\) be a real symmetric and positive definite matrix with an invertible~\(A\). Prove that the Schur complement \(S=C-BA^{-1}B^t\) is also symmetric and positive definite.
\item
Suppose~\(A\) is a Hermitian positive definite matrix split into \(A=C+C^*+D\) where~\(D\) is also Hermitian positive definite. Prove that \(B=C+\omega^{-1}D\) is invertible whenever \(0<\omega<2\). Consider the iteration \(x_{n+1}=x_n+B^{-1}(b-Ax_n)\), with any initial iterate \(x_0\). Prove that \(x_n\) converges to \(x=A^{-1}b\) whenever \(0<\omega<2\).
\item
If a square matrix~\(A\) can be block partitioned as \(A=\begin{bmatrix}0&M\\N&0\end{bmatrix}\), prove that \(\sigma(A)=\{z\in\mathbb{C}:z^2\in\sigma(MN)\cup\sigma(NM)\}\).
\item
Suppose~\(A\) is a Hermitian positive definite matrix with eigenvalues 
\[
0<\lambda_0\ll\lambda_1\leq\lambda_2\leq\cdots\leq\lambda_N.
\]
 Let \(\kappa=\lambda_N/\lambda_0\) and \(\kappa'=\lambda_N/\lambda_1\). Let \(e_n\) denote the error in the~\(n\)th step of the conjugate gradient algorithm applied to \(Ax=b\), and \(\lVert y\rVert_A=\langle Ay,y\rangle^{1/2}\). Prove that 
\[
\lVert e_n\rVert_A\leq 2\kappa\left(\frac{\sqrt{\kappa'}-1}{\sqrt{\kappa'}+1}\right)^{n-1}\lVert e_0\rVert_A.
\]
\item
Let~\(n\) be a positive integer, \(h=1/n\), and consider the grid of points \((ih,jh)\) for \(i,j=0,1,\ldots,n\). Let~\(A\) be the finite difference operator with the “5-point stencil” discretizing the Laplace operator \(-\Delta\), with zero Dirichlet boundary conditions on this grid. Describe it. What are the eigenvalues of~\(A\)? Prove that the spectral condition number of~\(A\) is \(\kappa(A)=\cot^2(\pi h/2)\).
\item
Let \(x_0,x_1,\ldots,x_n\) be distinct points in a finite interval \([a,b]\) and \(f\in C^1[a,b]\). Show that for any given \(\varepsilon>0\) there exists a polynomial~\(p\) such that \(\lVert f-p\rVert_\infty<\varepsilon\) and \(p(x_i)=f(x_i)\) for all \(i=0,1,\ldots,n\) (where \(\lVert\cdot\rVert_\infty\) denotes the \(L^\infty(a,b)\)-norm).
\item
Let \(w(x)\) be a positive integrable function on \([a,b]\). Consider the quadrature 
\[
\int_a^b f(x)w(x)\,dx\approx\sum_{k=0}^n\left(A_kf(x_k)+B_kf'(x_k)+C_kf''(x_k)\right)
\]
 for any~\(f\) with continuous first and second derivatives (\(f'\) and \(f''\), respectively). Give conditions on \(A_k,B_k,C_k\) and \(x_k\) so that the quadrature has precision \(4n+3\).
\item
Given that the Newton iteration for finding a root~\(r\) of~\(f\) converges, \(f\in C^2(\mathbb{R})\), and \(f(r)=f'(r)=0\neq f''(r)\), prove that the convergence cannot be quadratic. Suggest a modification that restores quadratic convergence for smooth~\(f\).
\item
Let \(x_0,x_1,\ldots,x_n\) be distinct numbers and for each~\(i\), let \(\ell_i\) denote the product of all \((x_i-x_j)^{-1}\) for all \(j\neq i\), \(j=0,\ldots,n\). Express the number \(q_f=f(x_0)\ell_0+f(x_1)\ell_1+\cdots+f(x_n)\ell_n\) as a divided difference of~\(f\). Then show that \(q_f=0\) whenever~\(f\) is a polynomial of degree \(n-1\).
\end{examproblems}
\end{document}
